\documentclass{article}

\usepackage{amsmath, amssymb}
\usepackage{geometry}
\geometry{margin=1in}

\title{Mathematics of the SNTRUP761-X25519-SHA512 Hybrid Key Exchange}
\author{}
\date{}

\begin{document}

\maketitle

\section{Overview}

The hybrid key exchange combines:

\begin{itemize}
\item X25519 elliptic-curve Diffie--Hellman
\item SNTRUP761 post-quantum lattice KEM
\item SHA-512 hashing
\end{itemize}

Two shared secrets are computed independently and then combined using a cryptographic hash.

\section{X25519 Key Exchange}

The protocol operates over the prime field

\[
\mathbb{F}_p, \quad p = 2^{255} - 19
\]

Let \(G\) be the base point on Curve25519.

\subsection{Key Generation}

Each party chooses a private scalar

\[
a \in \{0,1\}^{255}
\]

and computes the public key

\[
A = aG
\]

Similarly, the other party generates

\[
b \rightarrow B = bG
\]

\subsection{Shared Secret}

Both parties compute

\[
K_{x25519} = aB = abG
\]

which is identical on both sides.

\section{SNTRUP761 Post-Quantum Key Encapsulation}

Operations occur in the polynomial ring

\[
R = \mathbb{Z}_q[x] / (x^{761} - x - 1)
\]

with modulus

\[
q = 4591
\]

Polynomials have degree less than 761.

\subsection{Key Generation}

Choose small polynomials

\[
f, g \in \{-1,0,1\}^{761}
\]

Compute the modular inverse

\[
f^{-1} \pmod q
\]

The public key is

\[
h = g f^{-1} \pmod q
\]

The private key is

\[
(f,g)
\]

\subsection{Encapsulation}

Choose random small polynomial

\[
r \in \{-1,0,1\}^{761}
\]

Compute ciphertext

\[
c = r h + m \pmod q
\]

where \(m\) is an encoded message polynomial.

The shared secret is

\[
K_{sntrup} = H(m)
\]

\subsection{Decapsulation}

Receiver computes

\[
a = f c \pmod q
\]

Recover message

\[
m = a \pmod 3
\]

Then recompute the shared secret

\[
K_{sntrup} = H(m)
\]

\section{Hybrid Combination}

The two shared secrets are concatenated and hashed.

Let

\[
K_1 = K_{sntrup}
\]

\[
K_2 = K_{x25519}
\]

The final shared key is

\[
K = \text{SHA512}(K_1 \parallel K_2)
\]

\section{Handshake Summary}

\subsection{Client}

\begin{enumerate}
\item Generate X25519 key pair \(a, A\)
\item Generate SNTRUP761 key pair
\item Send \( (A, pk_{sntrup}) \)
\end{enumerate}

\subsection{Server}

\begin{enumerate}
\item Generate \(b, B\)
\item Encapsulate SNTRUP producing ciphertext \(c\)
\item Compute \(K_{x25519} = bA\)
\item Send \( (B, c) \)
\end{enumerate}

\subsection{Client}

\begin{enumerate}
\item Compute \(K_{x25519} = aB\)
\item Decapsulate SNTRUP to obtain \(K_{sntrup}\)
\end{enumerate}

Both parties compute

\[
K = SHA512(K_{sntrup} \parallel K_{x25519})
\]

\section{Security Intuition}

Security relies on two independent assumptions:

\begin{itemize}
\item Elliptic-curve discrete logarithm hardness (X25519)
\item Lattice-based hardness (SNTRUP761)
\end{itemize}

The hybrid construction remains secure if at least one of these primitives remains secure.

\end{document}

